% Continuous -> discrete connection. The approximation-under-discretization
% result is Lemma~\ref{lem:approx-disc} in Appendix~\ref{app:nuisance-rate};
% this section gives the implementation-facing narrative and the required
% refinement rate. (Superseded 2026-07: previously an isotropic-Holder
% duplicate of lem:approx-disc; now points to the PSHAB version.)

\section{From continuous to discrete: implementation via discretization}
\label{sec:discretization}

The convergence theory (Appendix~\ref{app:nuisance-rate}) is stated for a
target $\eta_0$ in the piecewise sparse anisotropic Besov class on $[0,1]^d$,
but optimal-tree solvers---including our implementation \texttt{optimaltrees}
via GOSDT---fit trees to \emph{binary features}. This section explains how the
continuous-covariate theory transfers to the discretized features the algorithm
actually uses.

\paragraph{Discretization scheme.}
For each continuous covariate $X_j \in [0,1]$ we select $m_{n,j}$ thresholds
$0 < t_{j,1} < \cdots < t_{j,m_{n,j}} < 1$ and form the binary indicators
$X_j^{(k)} = \mathbf 1\{X_j \le t_{j,k}\}$, $k = 1, \dots, m_{n,j}$. This
replaces each continuous covariate with $m_{n,j}$ binary features, yielding
$\tilde d = \sum_j m_{n,j}$ binary covariates; the solver then operates on
$\{0,1\}^{\tilde d}$. A tree restricted to these features can split coordinate
$j$ only at the chosen thresholds, so the reachable partitions are exactly the
grid-restricted trees $\cT_L^{\mathrm{disc}}$ of Lemma~\ref{lem:approx-disc}.

\paragraph{Discretization does not degrade the rate.}
The central question is whether restricting splits to a fixed grid inflates the
approximation error. It does not, provided the grid is fine enough:
Lemma~\ref{lem:approx-disc} shows that the grid-restricted infimum matches the
continuous rate $L^{-2\alphabar/s}$ once the per-coordinate resolution satisfies
the refinement condition~\eqref{eq:grid-rate}, namely
$m_{n,j} \gtrsim n^{(\alpha_{\min}/\alpha_j)/(2\alphabar + s)}$. The argument is
immediate rather than delicate: the achievability construction underlying the
sparse anisotropic rate is \emph{itself} a dyadic-grid tree
\citep{xuStatisticalOptimalityOptimal2026}, so a sufficiently fine grid already
contains the optimal approximant. Rougher coordinates (smaller $\alpha_j$)
require more thresholds; inactive coordinates require none.

\paragraph{Binary covariates.}
When the covariates are already binary, $\cX = \{0,1\}^d$, no discretization or
smoothness is needed: a tree with enough leaves represents any function on the
active subcube exactly, and the rate is governed purely by the sparsity $s$
(Corollary~\ref{cor:binary-rate}). This is the cleanest and most directly
applicable regime, and it is the one our binary-covariate simulations occupy.

\paragraph{Implementation in \texttt{optimaltrees}.}
The package accepts continuous features and discretizes them automatically;
predictions on new data apply the stored thresholds, keeping training and
inference consistent. To realize the theory, the number of thresholds must grow
with $n$ at the polynomial rate of~\eqref{eq:grid-rate}. A fixed or only
logarithmically growing number of bins---as in the current default---meets the
\emph{binary} guarantee (Corollary~\ref{cor:binary-rate}) but not the continuous
one; achieving the continuous rate requires the refinement schedule above, which
we adopt as the recommended default for continuous covariates.

\paragraph{Interpretation.}
Discretization is standard in tree methods (CART, random forests, and GOSDT all
split at data-driven thresholds, equivalent to discretizing). Our analysis makes
the requirement precise: the sparse anisotropic approximation rate is preserved
under discretization exactly when the grid refines at the rate
of~\eqref{eq:grid-rate}, and the estimation side is controlled by the
$L$-leaf-tree entropy of Lemma~\ref{lem:complexity} on the binary feature space.
